% Chapter 11
\chapter{Conclusions} % Main chapter title
\label{Chapter11} % For referencing the chapter elsewhere, use

This chapter summarizes the main findings of the thesis. The work investigates tendon-driven mechanisms from a mechanism-centered perspective, with Lasso Gripper serving as the primary novel mechanism platform and a handheld laparoscopic instrument serving as a complementary application case under the same tendon-driven framework. The chapter highlights the main contributions, clarifies the role of each prototype, and outlines directions for future work.

The central contribution of the thesis is Lasso Gripper, a tendon-driven end-effector that uses a launch-and-retraction string loop as the grasping structure itself. Unlike conventional antipodal grippers, which require relatively precise fingertip placement and often apply concentrated contact forces, Lasso Gripper creates a projected capture region and then converts enclosure into holding force through controlled retraction. This mechanism demonstrates how adaptive capture, distributed contact, low distal complexity, and explicit tension management can be combined in a single robotic end-effector.

The study of Lasso Gripper also identifies several mechanism-level findings. String mass, surface texture, launch speed, loop length, approach direction, and retraction timing jointly determine whether the loop can form, reach the target, tighten, and remain stable during lifting. The qualitative demonstrations show that the mechanism can exploit necks, horns, frames, irregular surfaces, deformable objects, oversized targets, and moving targets under selected setup conditions. At the same time, the observed failures, including geometric miss, premature or late closure, slip after capture, entanglement, and excessive compression, show that loop-based capture requires careful control of geometry, timing, and tension.

The handheld laparoscopic instrument is presented as a complementary application-oriented system rather than as a device derived from Lasso Gripper. It examines the same broad tendon-driven variables, including proximal actuation, guided cable transmission, paired winding, tension preservation, sensing, and low-level control, under a more constrained medical-device setting. Its value in the thesis is therefore mainly engineering-oriented: it demonstrates how a compact cable-driven architecture can be packaged with multiple output degrees of freedom, quick-swap tool interfaces, input sensing, and conservative safety logic.

Together, the two prototypes show that tendon-driven mechanisms can be organized for different manipulation regimes. Lasso Gripper emphasizes open-loop adaptive capture through a free tensile loop, while the surgical instrument emphasizes repeatable motion through guided cable paths. Their shared theme is not a linear transfer from one device to the other, but the use of flexible tensile elements, proximal actuation, and tension management to reduce distal complexity and shape interaction behavior.

Several limitations remain. The experiments with Lasso Gripper are demonstration-focused and do not yet provide large-sample statistical validation, randomized pose testing, or full closed-loop perception and planning. The surgical instrument remains an engineering prototype and requires further validation of sterilizability, ergonomic performance, long-term cable durability, control repeatability, and clinical safety. Future work should therefore focus on quantitative testing, parameter mapping, improved tension sensing, perception-assisted timing, adaptive control for cable hysteresis and stretch, and application-specific safety validation.

In conclusion, this thesis establishes Lasso Gripper as a novel tendon-driven mechanism for adaptive robotic capture and uses a handheld laparoscopic instrument as a complementary application case within the same mechanism framework. The results support the broader argument that flexible tensile transmission, when paired with explicit tension management and proximal actuation, can enable useful manipulation behaviors across both open robotic grasping and constrained handheld instrumentation.
